\documentclass[11pt]{article}

\usepackage{maiacustom}

\begin{document}

\psettitle{Astronomy Question Bank}

These questions were carefully written or selected to prepare students for the astronomy team selection process in Brazil. Some questions are not original; those are marked with () before the statement. The question-bank template is the same one used by Professor Kevin Zhou. His work is invaluable, and many ideas in this set can be found in his handouts.

% \begin{psidea}{Idea title}{}
% Idea
% \end{psidea}

% \begin{psexample}{Example title}{}
% Example
% \end{psexample}

% \begin{pssolution*}{}{}
% Solution
% \end{pssolution*}

% \begin{psremark*}{Remark title}{}
% Remark
% \end{psremark*}

\section{Thermodynamics}
\pts{3}
\begin{pproblem}
    A galaxy contains on the order of \(10^{10}\) stars. Because of this immense number, we can model a galaxy as a cloud of ideal gas, in which each star is equivalent to a particle of the gas.
    
    The goal of this question is to use this theoretical model to study some properties of galaxies. To that end, we make the following assumptions:

    \begin{enumerate}[label=\roman*)]
        \item The galaxy is spherical and in hydrostatic equilibrium.
        \item The mass density of the galaxy is constant and equal to \(\rho\).
        \item The stellar masses are small enough that interstellar interactions can be neglected.
    \end{enumerate}
    
    \begin{alternativas}
        \item Considering an ideal-gas system, find an expression for the pressure in terms of the density \(\rho\), the temperature \(T\), the mass of each particle \(\mu\), and physical constants.
        
        \item In our theoretical model it does not make sense to think in terms of temperature, so we need a substitute for it. Using the equipartition theorem, find an expression for \(T(r)\) and \(P(r)\).
    \end{alternativas}
\begin{pssolution*}{}{}
    \begin{alternativas}
        \item Starting from the Clapeyron equation,
        \[PV = Nk_BT\]

        Multiply both sides by \(\mu\),

        \[P\mu V = N\mu k_B T\]

        \[P\mu = \frac{N\mu}{V}k_BT\]

        Note that \(N\mu\) equals the total mass, so

        \[\rho = \frac{N\mu}{V}\]
        
        With this we obtain

        \[\boxed{P\mu = \rho k_B T}\]

        This expression will be used quite often in the thermodynamics part of this set.

        \item By the equipartition theorem,
        
        \[\frac{\mu v^2}{2} = \frac{3k_B T}{2}\]

        Hence
    
        \[T = \frac{\mu v^2}{3k_B}\]

        Since the stars have mass and speed, this is an acceptable substitution. We can compute \(v\) from the vis-viva equation, assuming circular orbits,

        \[v^2 = \frac{GM}{r^2} = \frac{4\pi G\rho r}{3}\]

        Thus

        \[\boxed{T(r) = \frac{4\pi G \rho \mu r}{9k_B}}\]

        For \(P(r)\), substitute into the formula:

        \[\boxed{P(r) = \frac{\rho k_B T(r)}{\mu} = \frac{4\pi G \rho^2 r}{9}}\]
    \end{alternativas}
    
\end{pssolution*}
\end{pproblem}

\pts{2}
\begin{pproblem}
    In this question we study the adiabatic index of stars. Assume that the exterior of a star is a vacuum at temperature \(T=0\).
    \begin{alternativas}
        \item All stars are bodies in hydrostatic equilibrium. Given that, what is the pressure at the surface of a star of mass \(M\) and radius \(R\)?
        \item Now suppose that the star expands by \(\delta R\). How does the pressure change? If needed, use \((1+x)^n\approx 1+nx\).
        \item Now conclude what minimum value of \(\gamma\), \(\gamma_{min}\), the star must have in order to remain gravitationally bound (assume that it expands adiabatically).
    \end{alternativas}

\begin{pssolution*}{}{}
    \begin{alternativas}
        \item Using \(P = F/A\), we have
        \[P(R) = \frac{GM^2}{4\pi R^4}\]

        \item Thus
        \[P(R+\delta R)  = \frac{GM^2}{4\pi}(R+\delta R)^{-4} = \frac{GM^2}{4\pi R^4}\left(1+\frac{\delta R}{R}\right)^{-4}\]

        Using the approximation given in the statement,

        \[\boxed{P(R+\delta R) = \frac{GM^2}{4\pi R^4}\left(1-\frac{4\delta R}{R}\right) = P\left(1-\frac{4\delta R}{R}\right)}\]
    
        \item In an adiabatic expansion, \(PV^\gamma\) is constant. Since \(V\propto R^3\), it follows that \(PR^{3\gamma}\) is constant. Thus
        
        \[PR^{3\gamma} = (P+\delta P)(R + \delta R)^{3\gamma}\]

        Using the same reasoning as in the previous item,
        \[PR^{3\gamma} = (P+\delta P)R^{3\gamma}\left(1 + \frac{\delta R}{R}\right)^{3\gamma}\]

        \[P = (P+ \delta P)\left(1+\frac{3\gamma\delta R}{R}\right)\]

        Substituting \(P+\delta P\),

        \[1 = \left(1-\frac{4\delta R}{R}\right)\left(1+\frac{3\gamma\delta R}{R}\right)\]

        Hence

        \[\left(1-\frac{4\delta R}{R}\right)^{-1} = 1+\frac{3\gamma\delta R}{R}\]

        Using the binomial approximation again on the left-hand side and solving for \(\gamma\), we obtain

        \[\boxed{\gamma = \frac{4}{3}}\]
    \end{alternativas}
    
\end{pssolution*}
\end{pproblem}

\pts{4}
\begin{pproblem}
    There are several models for the atmosphere of our planet. Let us explore them and find the physical effects of each one.
    \begin{alternativas}
        \item First, consider the isothermal model (\(T=\) const). Given that each air particle has mass \(\mu\) and that the atmosphere has \(P(0) = P_0\), find a formula for the pressure as a function of height, \(P(h)\).
        \item A more realistic model of the atmosphere is in fact adiabatic, since air is a very poor conductor of heat. Given that the air has Poisson coefficient \(\gamma\), find an expression for \(P(h)\) in the adiabatic model. Assume that at sea level the pressure and the temperature are \(P(0)\) and \(T(0)\).
        \item Find an expression for \(dT/dh\) for the previous model and estimate its value. Is the result consistent with reality?
    \end{alternativas}

\begin{pssolution*}{}{}
    \begin{alternativas}
        \item From the relation
        \[P\mu = \rho k_BT\]

        Thus, since \(T = \mathrm{const}\),

        \[\frac{P(h)}{P(0)} = \frac{\rho(h)}{\rho(0)}\]

        The pressure gradient is given by Stevin's law,

        \[\frac{dP(h)}{dh} = -g\rho(h)\]

        Substituting \(\rho(h)\),

        \[\frac{dP(h)}{dh} = \frac{-gP(h)\rho(0)}{P(0)}\]

        Separating the terms,

        \[\frac{dP(h)}{P(h)} = \frac{-g\rho(0)}{P(0)}dh\]

        Integrating both sides,

        \[\int_0^h\frac{dP(h)}{P(h)} = \ln\left(\frac{P(h)}{P(0)}\right)\]

        \[\int_0^hdh = h\]

        Thus

        \[P(h) = P(0)e^{-\frac{g\rho(0)}{P(0)}h}\]

        Finally, we eliminate the term \(\rho(0)\) using the Clapeyron relation,

        \[P(0)\mu = \rho(0)k_B T \rightarrow \rho(0) = \frac{P(0)\mu}{k_BT}\]

        Substituting,

        \[\boxed{P(h) = P(0)e^{-\frac{\mu g h}{k_BT}}}\]

        \item In the adiabatic model we have
        \[PV^\gamma = \mathrm{const}\]

        Since \(V\propto \rho^{-1}\),

        \[P\rho^{-\gamma} = \mathrm{const}\]

        That is,

        \[\rho(h) = \left(\frac{P(0)}{P(h)}\right)^{-1/\gamma}\rho(0) = \left(\frac{P(h)}{P(0)}\right)^{1/\gamma}\rho(0)\]

        Substituting into the expression for the pressure gradient,

        \[\frac{dP(h)}{dh} = -g\rho(h) = -\left(\frac{P(h)}{P(0)}\right)^{1/\gamma}\rho(0)g\]

        Separating the terms and integrating,

        \[\int_0^h P(h)^{-1/\gamma}dP(h) = -P(0)^{-1/\gamma}\rho(0)g\int_0^h dh\]

        \[\frac{\gamma}{\gamma-1}(P(h)^{\frac{\gamma-1}{\gamma}}-P(0)^{\frac{\gamma-1}{\gamma}}) = -\frac{\mu g z}{k_B T(0)}P(0)^{\frac{\gamma-1}{\gamma}}\]

        Here we use

        \[P(0)\mu = \rho(0)k_B T(0)\]

        Solving for \(P(h)\),
    
        \[\boxed{P(h) = P(0)\left(1-\frac{\gamma-1}{\gamma}\frac{\mu g h}{k_B T(0)}\right)^{\frac{\gamma}{\gamma-1}}}\]
       
        \item For adiabatic processes,
        
        \[PV^\gamma = \mathrm{const}\]

        \[V\propto \frac{T}{P}\]

        Hence

        \[P^{1-\gamma} T^\gamma = \mathrm{const}\]

        \[T P^\frac{1-\gamma}{\gamma} = \mathrm{const}\]
        
        Thus

        \[T = P(0)^\frac{1-\gamma}{\gamma}T(0)P^{\frac{\gamma-1}{\gamma}}\]

        Differentiating,

        \[\frac{dT}{dh} = P(0)^\frac{1-\gamma}{\gamma}T(0) \frac{\gamma-1}{\gamma P^{1/\gamma}}\frac{dP}{dh}\]

        From the previous item we have the expression for \(P\) as a function of \(h\). Carrying out the derivative,

        \[\frac{dP}{dh} = P(0)\frac{d}{dh}\left(1-\frac{\gamma-1}{\gamma}\frac{\mu g h}{k_B T(0)}\right)^{\frac{\gamma}{\gamma-1}}\]

        Using the chain rule,

        \[\frac{dP}{dh} = \frac{\gamma}{\gamma-1}P(0)\left(1-\frac{\gamma-1}{\gamma}\frac{\mu g h}{k_B T(0)}\right)^{\frac{1}{\gamma-1}}\left(-\frac{\gamma-1}{\gamma}\frac{\mu g}{k_B T(0)}\right)\]

        Simplifying,

        \[\frac{dP}{dh} = -P(0)\left(1-\frac{\gamma-1}{\gamma}\frac{\mu g h}{k_B T(0)}\right)^{\frac{1}{\gamma-1}}\frac{\mu g}{k_B T(0)}\]

        Substituting into the expression for \(dT/dh\),

        \[\frac{dT}{dh} = -P(0)^\frac{1-\gamma}{\gamma}T(0) \frac{\gamma-1}{\gamma P^{1/\gamma}}P(0)\left(1-\frac{\gamma-1}{\gamma}\frac{\mu g h}{k_B T(0)}\right)^{\frac{1}{\gamma-1}}\frac{\mu g}{k_B T(0)}\]

        \[\frac{dT}{dh} = -\frac{\gamma-1}{\gamma}P(0)^\frac{1}{\gamma} P^{-\frac{1}{\gamma}}\left(1-\frac{\gamma-1}{\gamma}\frac{\mu g h}{k_B T(0)}\right)^{\frac{1}{\gamma-1}}\frac{\mu g}{k_B}\]

        Substituting \(P\),

        \[\frac{dT}{dh} = -\frac{\gamma-1}{\gamma}P(0)^\frac{1}{\gamma} \left(P(0)\left(1-\frac{\gamma-1}{\gamma}\frac{\mu g h}{k_B T(0)}\right)^{\frac{\gamma}{\gamma-1}}\right)^{-\frac{1}{\gamma}}\left(1-\frac{\gamma-1}{\gamma}\frac{\mu g h}{k_B T(0)}\right)^{\frac{1}{\gamma-1}}\frac{\mu g}{k_B}\]

        Continuing to simplify, we obtain a neat expression,

        \[\frac{dT}{dh} = -\frac{\gamma-1}{\gamma}\frac{\mu g}{k_B}\]

        We can approximate air as a diatomic gas, \(\gamma = \frac{7}{5}\), of mass \(\mu\approx 30\)u. Carrying out the calculation, we obtain

        \[\boxed{\frac{dT}{dh} \approx - 10.03 \text{ K/km}}\]

        which is a value quite consistent with reality.
    \end{alternativas}
\end{pssolution*}
\end{pproblem}

\pts{5}
\begin{pproblem}
    Black holes are among the most fascinating bodies in the universe. In this question we study some of the thermodynamics associated with these objects. For this question we use natural units, i.e.\ \(c = G = \hbar = k_B = 1\). In this convention, the mass of the black hole is described by the equation

    \[dM = \frac{\kappa}{8\pi}dA + \Omega dL + \Phi dQ\]

    Here the quantities refer to the event horizon: \(\kappa\) is the gravitational acceleration at the event horizon, \(A\) its area, \(\Omega\) the angular velocity of the black hole, \(L\) its moment of inertia, \(\Phi\) the electric potential, and \(Q\) its charge.

    \begin{center}
    \textbf{Part 1: Basic Thermodynamics}   
    \end{center}
    
    \begin{alternativas}
        \item One of the fundamental concepts of thermodynamics is entropy. Using what you know about it, briefly explain the inequality
        \[\oint dS \geq 0\]

        \item Of the 3 factors that govern the mass of a black hole (\(dA, \ dL, \ dQ\)), only \(dA\) obeys the same inequality as the entropy. Why is this always true?
        
        For the next items, consider a black hole with no spin and no charge.

        \item Bekenstein and Hawking proved the so-called \textit{Bekenstein-Hawking entropy}, which relates the entropy to the area of the black hole (remember that we are using natural units, so some dimensions may not make sense). Bekenstein and Hawking found that for a black hole \(S \equiv \frac{A}{4}\). The equation we then have is
        
        \[dM = \frac{\kappa}{2\pi}d\left(\frac{A}{4}\right)\]

        By analogy between \(dM\) and some state function, find the temperature of the black hole as a function of \(\kappa\).

        An important term is still missing from the previous formula: the pressure associated with the dark-energy density, \(\Lambda\).

        \item The pressure due to dark energy is given by
        \[P = -\frac{\Lambda}{8\pi}\]

        where \(\Lambda\) is constant. This shows that \(VdP\) vanishes, so it may be added freely to the previous expression. We can therefore conclude that the mass of the black hole is in fact related to another thermodynamic potential. Which one is it?

        \item Note that the volume, \(V=V\), and the entropy, \(S = A/4\), are no longer independent for black holes. Assuming that the event horizon of the black hole is a sphere, find a relation between \(S\) and \(V\). This is further evidence that the internal energy, \(U = U(S, V)\), is not the best thermodynamic potential to work with.
    \end{alternativas}

    \centering
    \textbf{Part 2: Carnot Cycle for Black Holes}
    
    \begin{alternativas}
        \item The goal of this part is to build a theoretical model for a Carnot cycle inside black holes. First, however, prove an important result: for black holes, adiabatic and isochoric processes must be equivalent.
        
        \item Calculate the heat capacity at constant pressure of a black hole. Your result should be something bizarre.
        
        \item Use the fact that \(Q = T \Delta S\) along the isotherms, together with the results of the previous parts, to calculate the efficiency of a black-hole Carnot engine and confirm that you
        obtain the Carnot efficiency. Marvel at how much faster this calculation is than the
        typical derivation of the Carnot efficiency, and note that you have also inadvertently
        calculated the efficiency of the Stirling cycle.

        If you are interested in the subject, there is an interesting paper that discusses cycles in black holes specifically, available \href{https://arxiv.org/pdf/1404.5982}{here}.
    \end{alternativas}

    \centering
    \textbf{Part 3: Lifetime and Evaporation of Black Holes}
    
    \begin{alternativas}
        \item As calculated in part 1, black holes have a temperature. As a result, they emit radiation as described by the Stefan-Boltzmann law. Given that, find a relation between the lifetime of a black hole and its mass \(M\).
    \end{alternativas}

    \begin{pssolution*}{}{}
        \begin{center}
            \textbf{Part 1: Basic Thermodynamics}
        \end{center}
        \begin{alternativas}
            \item The inequality represents the second law of thermodynamics. Briefly, it shows that the change in entropy is always positive, that is, everything tends toward disorder.
            
            \item Note that no law of physics prevents a black hole from losing angular velocity, \(dL < 0\), or charge, \(dQ<0\). When we examine the area of the black hole, however, something interesting appears. Consider the event horizon, given by a sphere of radius \(R_{sch}\). Since this radius is defined as the distance at which an object moving at the speed of light can no longer escape the gravitational attraction of the black hole, we can obtain it from energy conservation,
            
            \[-\frac{GMm}{R_{sch}} + \frac{mc^2}{2} = 0 \rightarrow R_{sch} = \frac{2GM}{c^2}\]
            
            Hence the area of the event horizon is given by

            \[A = 4\pi R_{sch}^2 = \frac{8\pi G^2M^2}{c^4}\]

            In natural units,

            \[A = 8\pi M^2\]

            Thus

            \[dA = 16\pi MdM\]

            A black hole cannot lose mass, since nothing can escape the event horizon. Thus \(dM>0\) and consequently \(dA>0\). We can therefore conclude that it is also valid that
            
            \[\oint dA \geq 0\]

            for black holes.

            \item From the mass-energy relation,
            
            \[U = M \rightarrow dU = dM\]

            (Recall that we are using natural units, so \(Mc^2\) = \(M\).) Thus, using the first law of thermodynamics,

            \[dU = TdS - PdV\]

            Hence

            \[T = \left.\frac{dU}{dS}\right\vert_V = \left.\frac{dM}{dS}\right\vert_V\]

            Using the entropy formula given in the statement (analogous to the one obtained in the previous item),

            \[dM = \frac{\kappa}{2\pi}\left(\frac{A}{4}\right) = \frac{\kappa}{2\pi}dS\]

            Using the equation from the previous item, we see that \(\kappa = 1/4M\). Substituting into the formula for \(T\),

            \[\boxed{T = \frac{1}{8\pi M}}\]

            \item Using the formula obtained for \(T\),
            
            \[dM = TdS\]

            Adding the vanishing term \(VdP\), we have

            \[dM = TdS + VdP\]

            Note that this is identical to the enthalpy, \(H\):

            \[H = U + PV \rightarrow dH = (TdS - PdV) +  (PdV - VdP) = TdS - VdP\]

            We can therefore say that the mass is equivalent to the enthalpy.

            \item The volume of a sphere is given by
            
            \[V = \frac{4\pi R^3}{3}\]

            while the area is given by \(A = 4\pi R^2\), so \(R = \sqrt\frac{A}{4\pi}\equiv \sqrt\frac{S}{\pi}\). Substituting into the volume formula,
            
            \[V = \frac{4\pi}{3}\left(\frac{S}{\pi}\right)^{3/2}\]

            Simplifying,

            \[\boxed{V = \frac{4}{3\sqrt{\pi}}S^{3/2}}\]

            That is, the volume and the entropy depend on each other. Since the internal energy is a function of the volume \textbf{and} the entropy, it does not make sense to say that it equals the mass of the black hole.
        \end{alternativas}
    
        \begin{center}
            \textbf{Part 2: Carnot Cycle for Black Holes}
        \end{center}

        \begin{alternativas}
            \item Using the relation from the previous item,
            
            \[dV = \frac{2}{\pi}S dS\]

            That is, an adiabatic cycle (\(dS = 0\)) is equivalent to an isochoric cycle (\(dV = 0\)).

            \item To calculate the heat capacity at constant pressure,
            
            \[C_P = \left.\frac{\partial Q}{\partial T}\right\vert_P = T\left.\frac{\partial S}{\partial T}\right\vert_P = \frac{dM}{dT}\] 
            
            Using the formula for the temperature,

            \[T = \frac{1}{8\pi M} \rightarrow M = \frac{1}{8\pi T}\]
            \[\boxed{C_P = \frac{dM}{dT} = -\frac{1}{8\pi T^2}}\]

            \item The efficiency of a cycle is given by

            \[\eta = \frac{W}{Q_H} = \frac{Q_C-Q_H}{Q_H} = 1-\frac{Q_C}{Q_H}\]

            where \(Q_H\) is the heat that enters the cycle during the hot phase and \(Q_C\) is the amount that enters it during the cold phase. Define \(T_H\) and \(T_C\) as the hot and cold temperatures, respectively. Suppose the cycle operates with \(1-2\) at the hot temperature and \(3-4\) at the cold temperature. Then

            \[Q_H = T_H\Delta S_{1\rightarrow 2}\]

            \[Q_C = T_C\Delta S_{3\rightarrow 4}\]

            Using the expression from item \(1.e\), we have

            \[S = \frac{3\sqrt{\pi}}{4}V^{2/3}\]

            Thus

            \[\Delta S_{1\rightarrow 2} = \frac{3\sqrt\pi}{4}\left(V^{2/3}_2 - V^{2/3}_1\right)\]
            \[\Delta S_{3\rightarrow 4} = \frac{3\sqrt\pi}{4}\left(V^{2/3}_4 - V^{2/3}_3\right)\]

            Hence the efficiency can be written as

            \[\eta = 1 - \frac{T_C\left(V^{2/3}_4 - V^{2/3}_3\right)}{T_H \left(V^{2/3}_2 - V^{2/3}_1\right)}\]

            But since, for black holes, adiabatic and isochoric processes are identical, \(V_1=V_3\) and \(V_2=V_4\), so

            \[\boxed{\eta = 1 -\frac{T_C}{T_H}}\]

            which is the Carnot efficiency!

            \begin{center}
                \textbf{Part 3: Lifetime and Evaporation of Black Holes}
            \end{center}
        \end{alternativas}
        \begin{alternativas}
            \item The Stefan-Boltzmann law says that the luminosity (power) is given by
            \[L = -\frac{dE}{dt} = - \frac{dM}{dt} = A\sigma T^4 \]

            Note that \(A\propto M^2\) and \(T\propto 1/M\), so

            \[\frac{dM}{dt} \propto M^2\left(\frac{1}{M}\right)^4 = \frac{1}{M^2}\]

            Separating the variables,

            \[M^2dM \propto dt\]

            Finally, integrating, we obtain

            \[\boxed{t \propto M^3}\]
        \end{alternativas}
    \end{pssolution*}


\end{pproblem}

\pts{3}
\begin{pproblem}
    Consider a rocket that uses an ideal diatomic gas as fuel. Its mechanism is quite simple: the gas starts in a chamber at temperature \(T_1\) with cross-sectional area \(A_1\); the gas then flows adiabatically and is expelled through an opening of area \(A_2\), at pressure \(P_2\) and temperature \(T_2<T_1\). Assuming that the flow is continuous, determine the thrust felt by the rocket.
\begin{pssolution*}{}{}
    Since the process is adiabatic, we have

    \[PV^\gamma \propto P\left(\frac{T}{P}\right)^\gamma = \mathrm{const}\]

    Since the gas is diatomic, \(C_P = 7/2 R\) and \(C_V = 5/2 R\), and by definition \(\gamma = C_P/C_V = 7/5\).

    Thus we have

    \[P_1\left(\frac{T_1}{P_1}\right)^\gamma = P_2\left(\frac{T_2}{P_2}\right)^\gamma\]

    Substituting the value of \(\gamma\),

    \[\frac{P_1}{P_2} = \left(\frac{T_1}{T_2}\right)^{7/2}\]

    Let \(v_1\) and \(v_2\) be the speed of the gas at the given stages. Since the flow is continuous, we must have

    \[\rho_1v_1A_1 = \rho_2v_2A_2\]

    And by the ideal gas law,

    \[P\mu = \rho R T \rightarrow \rho \propto \frac{P}{T}\]

    Substituting into the previous expression,

    \[\frac{P_1v_1A_1}{T_1}=\frac{P_2A_2A_2}{T_2}\]

    Combining this with the first equation, we have
    
    \[\frac{v_1}{v_2} = \frac{A_2}{A_1}\left(\frac{T_2}{T_1}\right)^{5/2}\]

    Using Bernoulli's equation for gases,

    \[\frac{1}{2}\mu v_1^2 + C_P RT_1 = \frac{1}{2}\mu v_2^2 + C_P T_2\]

    Solving for \(v_2\),

    \[v_2^2 = \frac{7R(T_1-T_2)}{\mu(1-(A_1/A_2)^2(T_2/T_1)^5)}\]

    Finally, using Newton's second law,

    \[F = \frac{dp}{dt} = \frac{dm}{dt}v_2 = \rho_2A_2 v_2^2 = \boxed{\frac{7P_2A_2(T_1-T_2)}{T_2(1-(A_1/A_2)^2(T_2/T_1)^5)}} \]

\end{pssolution*}

\end{pproblem}


\pts{5} 
\begin{pproblem} (Adapted from Iran Problem Set)
    This problem aims to calculate the boiling point of liquids.

Particles in a liquid move at different speeds depending on the temperature, and some of these particles can escape the intermolecular forces and Earth's gravity (which will be neglected in this problem), leaving the surface of the liquid. These particles transfer their momentum, creating pressure when they collide with the surrounding environment. This pressure is known as the vapor pressure of the liquid. The boiling point is the temperature at which the vapor pressure equals the atmospheric pressure around the liquid.

\begin{alternativas}
\item Using the Maxwell-Boltzmann distribution, find a relation for the root-mean-square speed $v_{\text{rms}}$.

The Maxwell-Boltzmann distribution is
\begin{equation}
    n(v) \, dv = n \left(\frac{m}{2\pi k_B T}\right)^{3/2} e^{-\frac{mv^2}{2k_BT}} \, 4\pi v^2 dv
\end{equation}

The root-mean-square speed is given by
\begin{equation}
    v_{\text{rms}}^2 = \frac{1}{N} \sum_{i=1}^N v_i^2 = \frac{1}{n} \int_0^\infty v^2 n(v) dv
\end{equation}

\item Suppose that the speed of the particle under study equals $v_{\text{rms}}$. Also suppose that Earth's atmosphere is 80\% nitrogen and 20\% oxygen, and that the liquid under study is water.

Calculate the distance the escaped particle travels in the atmosphere before colliding, known as the mean free path. Express this distance in terms of the number density of the atmosphere and the geometric cross section of the particles.

\item Calculate the rate of change of the momentum of an escaping particle. Divide the change in momentum by the time interval of the process. Assume that the particles of the liquid lose all of their momentum when they collide with air molecules.

The mean time until the next collision is the mean free path divided by the speed of the particle.

Given that, over this time interval, a momentum equal to the momentum of the liquid particle is transferred to the air molecule, use Newton's second law to calculate the force exerted by the liquid particle on the air molecule.

\item Pressure is the force exerted on a surface. Find an expression for the vapor pressure of a liquid. You may leave your answer in terms of the cross sections of water and of air, \(S_w\) and \(S_a\), respectively.

\item Using the hydrostatic-equilibrium relation and assuming constant gravitational acceleration, constant air density, and zero pressure in the upper layers of the atmosphere, find a relation for the pressure near Earth's surface. Express this relation in terms of the air density, the gravitational acceleration, and the thickness of the atmosphere.

\item Impose the condition required for boiling by equating the atmospheric pressure near Earth's surface (obtained above) to the vapor pressure. Simplify the result until you obtain
\begin{equation}
    T = \frac{mhg S_w}{3k_B S_a}
\end{equation}

where $m$ is the mean mass of the air molecules, $g$ is the gravitational acceleration, $h$ is the thickness of the atmosphere, and the other parameters were described in the previous parts.

\item Determine the mean mass of the air molecules for the composition mentioned at the beginning of the problem.

\item Use the Bohr radius to estimate the ratio of the cross sections. Suppose an electron in a circular orbit around a proton, with the Coulomb force as the dominant force. Using Niels Bohr's assumption $L = n\hbar$, find the distance from the electron to the nucleus in terms of physical constants, $n$, and the atomic number $Z$.

\item Calculate the cross section of the air atoms and of the liquid atoms. Assume that each liquid molecule consists of two hydrogen atoms and one oxygen atom, and that the air particles consist of two oxygen atoms and two nitrogen atoms. Find the ratio of the cross sections $S_i/S_a$.

\item Take $g = 9.8 \, \text{m/s}^2$ and the height of the atmosphere to be $100 \, \text{km}$. Determine the boiling point of water.

\end{alternativas}

\begin{pssolution*}{}{}
    \begin{alternativas}
    \item Defining the variables,
    \[C = 4\pi n \left(\frac{m}{2\pi k_B T}\right)^{3/2}, \ \ \beta = \frac{1}{k_BT}\]

    we can write

    \[n(v)dv = C e^{-\frac{\beta mv^2}{2}}v^2dv\]

    Using the expression for \(v_{rms}\) given in the question,

    \[v_{rms}^2 = \frac{C}{n}\int_0^\infty  e^{-\frac{\beta mv^2}{2}}v^4 dv\]

    To evaluate the integral, we have

    \[\int_0^{\infty} e^{-av^2}v^4 dv = \frac{d^2}{da^2}\int_0^{\infty} e^{-av^2}dv\] 

    where

    \[a = \frac{\beta m}{2}\]
    The integral on the right-hand side is well known and equals \(\frac{1}{2}\sqrt{\frac{\pi}{a}}\). That is,

    \[\int_0^{\infty} e^{-av^2}v^4 dv = \frac{\sqrt{\pi}}{2}\frac{d^2}{da^2}a^{-1/2} = \frac{3}{8}\sqrt{\frac{\pi}{a^5}}\]

    Thus

    \[v_{rms}^2 = \frac{3C}{8n}\sqrt{\frac{\pi}{a^5}}\]

    Substituting \(C\) and \(a\), we have

    \[v_{rms}^2 = \frac{3}{2}\pi  \left(\frac{\beta m}{2\pi}\right)^{3/2}\sqrt{\frac{2^5\pi}{\beta^5m^5}}\]
    
    Finally, simplifying,

    \[\boxed{v_{rms}^2 = \frac{3}{m\beta} = \frac{3k_B T}{m}}\]

    \item Consider \(P(t)\) the probability that the particle does \textbf{not} collide during a time \(t\). The probability that the particle does not collide during a time \(t\) and does not collide during the following time \(dt\) is \(P(t+dt) = P(t)P(dt)\) (multiplicative property), but we also have
    
    \[P(t+dt) \approx P + \frac{dP}{dt}dt\]

    In a time \(dt\) the particle sweeps out a volume \(dV = \sigma v dt\), where \(\sigma\) is the impact parameter. The probability that the particle collides in a time \(dt\) is then \(P'(dt) = ndV = n\sigma v dt\), so (recalling that \(P(dt)\) is the probability that the particle does \textbf{not} collide) \(P(dt) = 1-n\sigma vdt\).

    Equating the expressions for \(P+dt\), we have

    \[P + \frac{dP}{dt}dt = P(1-n\sigma v dt)\]

    Working with this expression,

    \[dP = -Pn\sigma vdv \rightarrow \frac{dP}{P} = -n\sigma v dt\]

    Integrating both sides and using \(P(0)=1\), we have

    \[P(t) = e^{-n\sigma vt}\]

    Then the probability that a particle survives for a time \(t\) and then collides in the following time \(dt\) is

    \[J(t)dt = P(t)(1-P(dt)) = e^{-n\sigma v t} n\sigma v dt\]

    Thus the mean time between collisions is

    \[\tau = \int_0^\infty t J(t)dt = \int_0^\infty e^{-n\sigma v t} n\sigma v dt = \frac{1}{n\sigma v}\]

    The mean free path is then given by

    \[\lambda = v_{radial}\tau = \frac{1}{\sqrt{2}n\sigma}\]

    To find \(\sigma\), define \(r_H\) as the radius of the water molecule, and similarly \(r_O\) and \(r_N\). Looking at the following diagram,

    \begin{figure}[H]
        \centering
        \includegraphics[width= 0.7\linewidth]{imagens/secaotransversal.png}
        \caption{Source: Concepts in Thermal Physics, Blundell and Blundell}
    \end{figure}

    For nitrogen, \(\sigma_N = \pi(r_H + r_N)^2\), and for oxygen, \(\sigma_O = \pi (r_H + r_O)^2\). Since the atmosphere consists of \(0.8 N\) and \(0.2 O\), it holds that

    \[\sigma \equiv S_a = 0.2\sigma_O + 0.8\sigma N\]

    Thus we have

    \[\boxed{\lambda = \frac{1}{\sqrt{2}n(0.8\sigma_0 + 0.2\sigma_0)}}\]

    \item Assuming that the particle loses all of its momentum when it collides with an air molecule, we have \(\delta p = mv_{rms}\). Each collision occurs in a time \(\delta t =\tau\), so
    \[F = \frac{\delta p }{\delta t} = \frac{mv_{rms}}{\tau} = nS_a v_{rms}^2 \  \boxed{ = 3nS_a k_BT}\]

    \item Using the definition of pressure,
    
    \[P = \frac{F}{\S_w} = \ \boxed{3n\frac{S_a}{S_w}k_B T} \]

    \item Using the hydrostatic-equilibrium equation,
    
    \[\frac{dP}{dr} = -\rho g\]

    Since \(\rho\) and \(g\) are constant,

    \[P = -\rho g \int_h^0dh= \ \boxed{\rho g h}\]

    where \(h\) is the thickness of the atmosphere.

    \item Equating the pressures,

    \[\rho g h = 3\frac{S_a}{S_w}nk_BT\]

    Solving for \(T\), we have

    \[T = \frac{\rho g h}{3nk_B}\frac{S_w}{S_a}\]

    Since \(n\) is the number density of particles and \(\rho\) is the mass density, we have \(\rho/n = m\), so

    \[\boxed{T= \frac{m g h}{3k_B}\frac{S_w}{S_a}}\]

    \item Taking a weighted average,
    
    \[\boxed{m = 0.8 m_N + 0.2 m_O}\]

    \item The force felt by the electron is given by
    
    \[F = \frac{Ze^2}{4\pi \epsilon_0r^2}\]

    Equating this to the centripetal force,

    \[\frac{Ze^2}{4\pi \epsilon_0r^2} = m_e \omega^2 r\]

    By the definition of angular momentum, we have

    \[L = m_er^2\omega \rightarrow \omega = \frac{L}{m_er^2} = \frac{n\hbar}{m_er^2}\]

    Substituting into the previous expression,

    \[\frac{Ze^2}{4\pi\epsilon_0 r^2} = m_e \left(\frac{n\hbar}{m_er^2}\right)^2 r\]

    Solving for \(r\),

    \[r = \frac{4\pi\epsilon_0 n^2\hbar ^2}{Z e^2 m_e}\]

    \item For water we have \(r_w = 2r_H + r_O\), and for air \(r_a = 2(r_N + r_O)\). Since \(S_i \propto r_i^2\), we have
    
    \[\frac{S_w}{S_a} = \frac{(2r_H + r_O)^2}{4(r_N + r_O)^2}\]

    Assuming that all the electrons are in the lowest level (\(n=1\)), that \(Z_H = 1, \ Z_N = 7, \ Z_O = 8\), and that \(r\propto 1/Z\), we have
    
    \[\boxed{\frac{S_w}{S_a}=\frac{(2/Z_H + 1/Z_O)^2}{4(1/Z_N + 1/Z_O)^2} \approx 15.73}\]

    \item Substituting into the formula,

    \[T = \frac{(2\cdot 14 + 2 \cdot 16) \cdot 1.67 \cdot 10^{-27} \cdot 10^{5} \cdot 15.73}{3 \cdot 1.38\cdot 10^{-23}} \ \boxed{\approx 380\text{ K}}\]

    which is an estimate quite consistent with reality, \(T = 373\) K.
\end{alternativas}
    

    

    
    
\end{pssolution*}
\end{pproblem}

\pts{3}
\begin{pproblem} (Adapted from Iran Problem Set)
    Dr. Shahram Abbassi is one of the most renowned Iranian scientists in the field of accretion disks. In one of his recent studies of the giant molecular cloud B32, he found a Sun-like star at the center of that cloud. According to his research, the cloud has a mass of $10^6 M_\odot$, a radius of $30 \, \text{pc}$, and a very high viscosity, so high that if the cloud were to collapse, the whole system would collapse with spherical symmetry.

    What matters most is to calculate the specific heat at constant volume ($C_v$) of this cloud. Using the data given and reasonable approximations, determine a bound on $C_v$ such that accretion is possible. Do these values depend on the mass and the radius of the cloud? What can we conclude from the result?
\begin{pssolution*}{}{}
    This bound is given by the virial theorem, when

    \[\langle U \rangle = -2 \langle K \rangle\]

    We can model the cloud as a gas, so that its kinetic energy is given by \(\frac{3}{2} N k_B T\). Its potential energy is the potential energy of a sphere, given by \(U = -3GM^2/5R\). Thus

    \[\frac{3GM^2}{5R} = 3\frac{M}{m_p}k_B T\]

    where \(m\) is the mass of a particle in the cloud.

    Thus

    \[T = \frac{GMm_p}{5k_B R}\]

    In terms of the specific heat at constant volume, the energy is given by

    \[M C_v T\]

    Returning to the equality,

    \[C_v = \frac{3GM}{5R T}\]

    Substituting \(T\),

    \[\boxed{C_v = \frac{k_B }{m} \approx 4.1 \cdot 10^3 \text{ J kg\(^{-1}\) K\(^{-1}\)}}\]

    Note that this value depends only on the mass of the particles that make up the gas cloud. For a hydrogen composition, the value is surprisingly close to that of water, \(C_{v, \text{water}} = 4.2\cdot 10^{3}\) J/kg K.
\end{pssolution*}
\end{pproblem}


\end{document}
